\documentclass[11pt]{article}

\usepackage[margin=0.75in]{geometry}
\usepackage[T1]{fontenc}
\usepackage[utf8]{inputenc}
\usepackage{lmodern}
\usepackage{microtype}
\usepackage{xcolor}
\usepackage{hyperref}
\usepackage{enumitem}
\usepackage{booktabs}
\usepackage{longtable}
\usepackage{array}
\usepackage{tabularx}
\usepackage{fancyhdr}
\usepackage{titlesec}
\usepackage[most]{tcolorbox}

\definecolor{GitHubBlack}{HTML}{24292F}
\definecolor{GitHubBlue}{HTML}{0969DA}
\definecolor{LightBlue}{HTML}{EFF6FF}
\definecolor{LightGray}{HTML}{F6F8FA}
\definecolor{RuleGray}{HTML}{D0D7DE}
\definecolor{WarnAmber}{HTML}{FFF8C5}

\hypersetup{
  colorlinks=true,
  linkcolor=GitHubBlue,
  urlcolor=GitHubBlue,
  citecolor=GitHubBlue,
  pdftitle={GitHub, GitHub Actions, and GHAS as a DevSecOps CI/CD Foundation},
  pdfauthor={Jordan Suber},
  pdfsubject={Citation-clean DevSecOps handout},
  pdfkeywords={GitHub, GitHub Actions, GHAS, GitHub Advanced Security, DevSecOps, CI/CD, CodeQL, secret scanning, Dependabot}
}

\setlength{\headheight}{14pt}
\pagestyle{fancy}
\fancyhf{}
\lhead{DevSecOps CI/CD Foundation}
\rhead{GitHub + GitHub Actions + GHAS}
\cfoot{\thepage}
\renewcommand{\headrulewidth}{0.4pt}

\titleformat{\section}{\large\bfseries\color{GitHubBlack}}{}{0em}{}[\titlerule]
\titleformat{\subsection}{\normalsize\bfseries\color{GitHubBlack}}{}{0em}{}
\setlist[itemize]{leftmargin=1.3em, itemsep=0.25em, topsep=0.25em}
\setlist[enumerate]{leftmargin=1.5em, itemsep=0.25em, topsep=0.25em}
\renewcommand{\arraystretch}{1.18}

\newtcolorbox{positionbox}{
  enhanced,
  colback=LightBlue,
  colframe=GitHubBlue,
  boxrule=0.8pt,
  arc=2mm,
  left=8pt,
  right=8pt,
  top=8pt,
  bottom=8pt
}

\newtcolorbox{cautionbox}{
  enhanced,
  colback=WarnAmber,
  colframe=RuleGray,
  boxrule=0.6pt,
  arc=2mm,
  left=8pt,
  right=8pt,
  top=8pt,
  bottom=8pt
}

\newcommand{\term}[1]{\textbf{#1}}
\newcommand{\artifact}[1]{\texttt{#1}}

\begin{document}

\begin{center}
  {\LARGE \textbf{GitHub, GitHub Actions, and GHAS}}\\[0.25em]
  {\large A Citation-Clean DevSecOps CI/CD Foundation Handout}\\[0.5em]
  {\small Prepared for enterprise AppSec, CI/CD, and platform engineering discussions}
\end{center}

\vspace{0.5em}

\begin{positionbox}
\textbf{Corrected position.} GitHub, GitHub Actions, and GitHub Advanced Security - now commonly discussed through GitHub Secret Protection and GitHub Code Security capabilities - provide a strong DevSecOps-ready CI/CD foundation. GitHub is the source-control and collaboration layer; GitHub Actions is the workflow automation layer for build, test, scan, and deployment pipelines; and GHAS/GitHub Code Security adds native AppSec and software supply-chain controls such as code scanning, secret scanning, Dependabot alerts, and dependency review. This stack is a foundation for enterprise delivery, but it does not by itself replace artifact management, release orchestration, ITSM/change control, runtime security, observability, or governance workflows.\cite{github-actions-understanding,github-ghas-about}
\end{positionbox}

\section{Executive Summary}

A precise enterprise framing is:

\begin{quote}
\textbf{GitHub is the source-control and collaboration platform; GitHub Actions is the workflow automation layer for CI/CD; and GitHub Advanced Security/GitHub Code Security is the native AppSec layer for code, secret, and dependency risk. Together, they provide a DevSecOps-ready CI/CD foundation that can be extended with deployment orchestration, artifact management, runtime security, ITSM, and governance tooling for a complete enterprise delivery platform.}
\end{quote}

GitHub Actions is officially described as a continuous integration and continuous delivery platform for automating build, test, and deployment pipelines.\cite{github-actions-understanding,github-actions-quickstart} GitHub Advanced Security documentation describes security capabilities grouped around GitHub Secret Protection and GitHub Code Security, including secret scanning, push protection, code scanning, premium Dependabot features, and dependency review.\cite{github-ghas-about}

\section{Component Responsibilities}

\begin{longtable}{>{\raggedright\arraybackslash}p{0.19\linewidth} >{\raggedright\arraybackslash}p{0.37\linewidth} >{\raggedright\arraybackslash}p{0.35\linewidth}}
\toprule
\textbf{Component} & \textbf{Primary Role} & \textbf{Typical Enterprise Use} \\
\midrule
\endhead
\term{GitHub} & Repository hosting, pull requests, code review, branch protection, CODEOWNERS, collaboration, and repository governance. & Establish the system of record for source code, review workflows, repository metadata, ownership, and development history. \\
\midrule
\term{GitHub Actions} & CI/CD workflow automation for repository-based jobs, including builds, tests, scans, packaging, and deployments.\cite{github-actions-understanding} & Automate pull-request validation, build pipelines, security scan invocation, artifact publication, release preparation, and environment-specific deployment workflows. \\
\midrule
\term{GitHub Advanced Security / GitHub Code Security} & Native GitHub security capabilities for identifying and managing code, secret, and dependency risk.\cite{github-ghas-about} & Provide CodeQL-based code scanning, secret scanning, Dependabot alerts, dependency review, and security alert governance within the developer workflow. \\
\bottomrule
\end{longtable}

\section{Security Capability Map}

\begin{longtable}{>{\raggedright\arraybackslash}p{0.23\linewidth} >{\raggedright\arraybackslash}p{0.43\linewidth} >{\raggedright\arraybackslash}p{0.25\linewidth}}
\toprule
\textbf{Capability} & \textbf{What It Does} & \textbf{Primary Risk Area} \\
\midrule
\endhead
\term{Code scanning with CodeQL} & Uses CodeQL to identify vulnerabilities and coding errors in code; results are displayed as code scanning alerts in GitHub.\cite{github-codeql-about} & SAST, secure coding, vulnerability discovery. \\
\midrule
\term{Secret scanning} & Detects secrets such as API keys and tokens in repositories and can support prevention through push protection.\cite{github-ghas-about,github-security-settings} & Credential leakage and exposed secrets. \\
\midrule
\term{Dependabot alerts} & Creates alerts when GitHub detects vulnerable dependencies in a repository dependency graph; alerts include affected file details, vulnerability severity, and fixed-version information when available.\cite{github-dependabot-alerts} & Software composition analysis and dependency vulnerability management. \\
\midrule
\term{Dependency review} & Reviews dependency changes in pull requests so teams can avoid introducing vulnerable dependencies rather than fixing them later.\cite{github-dependency-review} & Pull-request supply-chain risk control. \\
\midrule
\term{Code scanning integrations} & Allows externally produced code scanning results to be uploaded into GitHub and displayed as code scanning alerts.\cite{github-code-scanning-integration} & Unified alert visibility and integration with third-party scanners. \\
\bottomrule
\end{longtable}

\section{Corrected Architecture Interpretation}

\subsection{What the Stack Is Strong At}

\begin{itemize}
  \item \textbf{Developer-native security feedback:} Findings appear in the same platform where developers review and change code.
  \item \textbf{Pull-request quality gates:} GitHub Actions can run validation jobs on pull requests and enforce required status checks through repository governance.
  \item \textbf{Repeatable CI/CD workflows:} Workflows are versioned with the repository, making build, test, and scan logic auditable.
  \item \textbf{Native AppSec controls:} CodeQL, secret scanning, Dependabot, and dependency review address common application-security and supply-chain risk categories.
  \item \textbf{Extensibility:} GitHub Actions can call external tools, publish artifacts, request approvals, and integrate with enterprise systems.
\end{itemize}

\subsection{What It Does Not Fully Replace}

\begin{cautionbox}
\textbf{Important distinction.} GitHub Actions can perform deployments, but enterprise deployment orchestration is broader than running a deployment job. A complete enterprise platform commonly requires environment promotion policy, artifact repository controls, deployment approvals, change-management evidence, runtime policy enforcement, observability, incident response, and rollback governance.
\end{cautionbox}

\begin{longtable}{>{\raggedright\arraybackslash}p{0.28\linewidth} >{\raggedright\arraybackslash}p{0.62\linewidth}}
\toprule
\textbf{Area} & \textbf{Why an Extension Is Usually Needed} \\
\midrule
\endhead
\term{Artifact management} & Enterprises usually need immutable versioned artifacts, promotion stages, provenance, retention, and repository-level controls outside the source repository. \\
\midrule
\term{Deployment orchestration} & Multi-environment deployments often require progressive delivery, approvals, environment locking, release windows, service dependencies, and rollback patterns. \\
\midrule
\term{ITSM/change governance} & Regulated environments often need change records, approvals, evidence capture, assignment groups, and audit-ready traceability. \\
\midrule
\term{Runtime security} & Code and dependency checks do not replace container runtime policy, Kubernetes admission controls, workload identity controls, or production telemetry. \\
\midrule
\term{Enterprise reporting} & Security leaders usually need portfolio-level risk reporting, SLA tracking, exception handling, and ownership mapped to CMDB or service inventory data. \\
\bottomrule
\end{longtable}

\section{Recommended Enterprise DevSecOps Control Flow}

\begin{enumerate}
  \item \term{Developer opens a pull request.} Repository rules, CODEOWNERS, and branch protections establish review expectations.
  \item \term{GitHub Actions runs CI checks.} Build, test, linting, packaging, and security jobs execute as workflow jobs.
  \item \term{CodeQL/code scanning evaluates code risk.} CodeQL analysis produces code scanning alerts for vulnerabilities and coding errors.\cite{github-codeql-about}
  \item \term{Secret scanning checks for exposed credentials.} Secret scanning and push protection reduce credential-leak risk.\cite{github-ghas-about,github-security-settings}
  \item \term{Dependency review evaluates pull-request changes.} New or changed dependencies can be reviewed before risky components enter the codebase.\cite{github-dependency-review}
  \item \term{Dependabot alerts monitor known vulnerable dependencies.} Alerts support ongoing dependency remediation after vulnerable dependencies are detected.\cite{github-dependabot-alerts}
  \item \term{Approved artifacts are promoted.} Publish and promote signed, immutable artifacts through controlled repositories and environments.
  \item \term{Deployment proceeds through release governance.} Connect workflow evidence to deployment records, environment approvals, ITSM change records, and runtime controls.
\end{enumerate}

\section{Minimum Policy Recommendations}

\begin{itemize}
  \item Require pull-request checks for build, test, CodeQL/code scanning, secret scanning status, and dependency review before merge.
  \item Treat open critical/high security findings as explicit release-risk inputs rather than passive alerts.
  \item Define who can dismiss or close security alerts, what evidence is required, and when a historical alert must be reopened.
  \item Map repositories to applications, services, owners, and CMDB/application inventory records.
  \item Preserve evidence from CI/CD runs, security scan results, artifact publication, deployment approval, and production release events.
  \item Separate CI validation from deployment authorization: a successful pipeline run should not automatically imply production authorization in regulated environments.
\end{itemize}

\section{Source Quality Notes}

The original source list mixed official documentation with blog posts, vendor comparison pages, social posts, and generic explanatory material. For an enterprise handout, the primary references should be official GitHub documentation. Third-party articles can be useful for commentary, product comparison, or implementation examples, but they should not be the evidentiary basis for defining GitHub, GitHub Actions, GHAS, CodeQL, secret scanning, Dependabot, or dependency review.

\section{Clean Reference Set}

{\small\sloppy
\begin{thebibliography}{9}

\bibitem{github-actions-understanding}
GitHub Docs. \textit{Understanding GitHub Actions}. Available at: \url{https://docs.github.com/articles/getting-started-with-github-actions}.

\bibitem{github-actions-quickstart}
GitHub Docs. \textit{Quickstart for GitHub Actions}. Available at: \url{https://docs.github.com/en/actions/get-started/quickstart}.

\bibitem{github-ghas-about}
GitHub Docs. \textit{About GitHub Advanced Security}. Available at: \url{https://docs.github.com/en/get-started/learning-about-github/about-github-advanced-security}.

\bibitem{github-codeql-about}
GitHub Docs. \textit{About code scanning with CodeQL}. Available at: \url{https://docs.github.com/code-security/code-scanning/introduction-to-code-scanning/about-code-scanning-with-codeql}.

\bibitem{github-dependabot-alerts}
GitHub Docs. \textit{About Dependabot alerts}. Available at: \url{https://docs.github.com/code-security/dependabot/dependabot-alerts/about-dependabot-alerts}.

\bibitem{github-dependency-review}
GitHub Docs. \textit{About dependency review}. Available at: \url{https://docs.github.com/code-security/supply-chain-security/understanding-your-software-supply-chain/about-dependency-review}.

\bibitem{github-security-settings}
GitHub Docs. \textit{Managing security and analysis settings for your repository}. Available at: \url{https://docs.github.com/enterprise-cloud@latest/repositories/managing-your-repositorys-settings-and-features/enabling-features-for-your-repository/managing-security-and-analysis-settings-for-your-repository}.

\bibitem{github-code-scanning-integration}
GitHub Docs. \textit{About integration with code scanning}. Available at: \url{https://docs.github.com/en/code-security/concepts/code-scanning/about-integration-with-code-scanning}.

\end{thebibliography}
}

\end{document}
